\documentclass[10pt]{article}
\usepackage{titlesec}
\usepackage{geometry}
\geometry{verbose,tmargin=.9in,bmargin=.9in,lmargin=1.0in,rmargin=1.0in}
\usepackage{amsmath,amsfonts,amsthm,amssymb}
\usepackage{url}
\usepackage{color}
\usepackage[usenames,dvipsnames,svgnames,table]{xcolor}
\usepackage[colorlinks=true, linkcolor=red, urlcolor=blue, citecolor=gray]{hyperref}
\usepackage{float}
\usepackage{caption}
\usepackage{subcaption}
\usepackage{graphicx}
\usepackage{wrapfig}
\usepackage{booktabs}
\usepackage{longtable}
\usepackage{enumitem}
\usepackage{multicol}
\usepackage{etoolbox}

\definecolor{nyuDarkPurple}{HTML}{330662}
\definecolor{nyuOfficialPurple}{HTML}{57068c}

\newcommand{\spara}[1]{\vspace{.5em}\noindent {\large\sffamily\textcolor{nyuOfficialPurple}{#1}}}
\titleformat{\section}[hang]{\Large\sffamily\color{nyuDarkPurple}}{\thesection}{1em}{}
\titleformat{\subsection}[hang]{\large\sffamily\color{nyuDarkPurple}}{\thesection}{1em}{}
\titleformat{\subsubsection}[hang]{\normalsize\sffamily\color{gray}}{\thesection}{1em}{}

\usepackage{fancyhdr}
\pagestyle{fancy}
\lhead{\includegraphics[width=4cm]{tandon_long_color.eps}}
\rhead{\thepage}
\pagenumbering{gobble}

\setcounter{secnumdepth}{0}

% math commands
\DeclareMathOperator{\R}{\mathbb{R}}
\newcommand{\E}{\mathbb{E}}

\begin{document}
	
\begin{center}
	\normalsize
	New York University Tandon School of Engineering
	
	Computer Science and Engineering
	\medskip
	
	\large
	CS-GY 6763: Homework 1. 
	
	Due Thursday, Feb. 6th, 2024, 11:59pm ET.
	\medskip
	
	\normalsize 
	\noindent \emph{Collaboration is allowed on this problem set, but solutions must be written-up individually. Please list collaborators for each problem separately, or write ``No Collaborators" if you worked alone.}
	\medskip
	
	\noindent \emph{For just this first problem set, $10\%$ extra credit will be given if solutions are typewritten (using LaTeX, Markdown, or some other mathematical formatting program).}

\end{center} 

\subsection{Problem 1: Short(er) problems. (8 pts)} 
\begin{enumerate}
	\item (4pts) Consider inserting $m$ keys into a hash table of size  $n = 5 m^2$ using a uniformly random hash function. By the mark-and-recapture analysis from class, we know that the expected number of collisions in the table is $\frac{m\cdot (m-1)}{2\cdot 5 m^2} < 1/10$. So, by Markov's inequality, with probability $> 9/10$, the table has no collisions ($<1$ collisions). Thus, we can look up items from the table in worst-case $O(1)$ time. 
	
	Give an alternative proof of the fact that we have no collisions with $> 9/10$ probability in a table of size $cm^2$ for some sufficiently large constant $c$. Specifically, to have no collisions, we must have the following events all happen in sequence: the second item inserted into the hash table doesn't collide with an existing item, the third item inserted doesn't collide with an existing item, $\ldots$, the $m^\text{th}$ item inserted doesn't collide with an existing item. Analyze the probability these events all happen. \textbf{Hint:} You might want to use the fact that $\frac{1}{2e} \leq (1-\frac{1}{n})^n \leq \frac{1}{e}$ for any positive integer $n \geq 2$. 
	
	\item (4pts) Consider a random walk on a $d$-dimensional grid. At each step, the walk chooses one of the $d$-dimensions uniformly at random, and takes a step in that direction – up with probability $1/2$ and down with probability $1/2$. Assume that the walk starts at the origin, takes $n$ steps, and ends at position $x$, where $x$ is a $d$ dimensional vector with integer values. Show that $\E[\|x\|_2^2] = n$, where $\|x\|_2^2 = \sum_{i=1}^d x_i^2$ denotes the squared Euclidean norm of $x$.  
	
	I find this fact surprising. If a tourist starts at Washington Square Park and wonders around Manhattan (i.e., takes a two-dimensions random walk), in expectation they don't get any more lost than if they restrict their wondering to just up and down 5th Avenue (i.e., a one dimensional random walk.)
\end{enumerate}


\subsection{Problem 2: Why does Count-Min work so well in practice? (12 pts)}
We showed that Count-Min can estimate the frequency of any item in a stream of $n$ items up to additive error $\frac{1}{m} n$ using $O(m)$ space. In practice it is often observed that this bound is pessimistic: the algorithm performs better than expected. In this problem, you will establish one reason why. 

For any positive integer $m$, let $f_1, \ldots, f_{m}$ be the frequencies of the $m$ most frequent items in our stream. Let $C = n - \sum_{i=1}^{m} f_i$. In general, we can have that $C \ll n$. For example, it has been observed that up to $95\%$ of YouTube video views come from just $1\%$ of videos. 
Prove that using $O(m)$ space,  Count-Min actually returns an estimate $\tilde{f}$ to $f(v)$ for any item $v$ satisfying:
\begin{align*}
	f(v) \leq \tilde{f} \leq f(v) + \frac{1}{m} C
\end{align*}
with $9/10$ probability. This is strictly better than the $\frac{1}{m} n$ error bound shown in class. 


\subsection{Problem 3: Randomized methods for efficient disease identification group testing. (12 pts)}
One of the most important factors in controlling diseases like bird flu or, a few years ago, COVID-19, is testing. However, testing often requires processing in a lab, so can be expensive and slow. One way to make it cheaper is to test patients/livestock/etc. in \emph{groups}. The biological samples from multiple individuals (e.g., multiple nose swabs) are combined into a single test tube and tested for the disease all at once. If the test comes back negative, we know everyone in the group is negative. If the test comes back positive, we do not know which patients in the group actually have the disease, so further testing would be necessary. There's a trade-off here, but it turns out that, overall, group testing can save on the total number of tests run. 

\begin{enumerate}
	\item (6pts) Consider the following deterministic ``two-level'' testing scheme. We divide a population of $n$ individuals to be tested into $C$ groups of the same size. We then test each of these groups. For any group that comes back positive, we retest all members of the group individually. Show that there is a choice for $C$ such that, if $k$ individuals in the population have a disease (would test positive), we can find all of those individuals with $\leq 2\sqrt{nk}$ tests. You can assume $k$ is known in advance (often it can be estimated accurately from the positive rate of prior tests). This is already an improvement on the naive $n$ tests when $k < 25\% \cdot n$.
	
	\item (6pts) We can use randomness to do better. Consider the following scheme: Collect $q = O(\log n)$ biological samples from each individual (in reality, divide one sample into $q$ parts). Then, repeat the following process $q$ times: randomly partition our set of $n$ individuals into $C$ groups, and test each group in aggregate. Once this process is complete, report that an individual ``is positive'' if the group they were part of tested positive all $q$ times. Report that an individual ``is negative'' if \emph{any} of the groups they were part of tested negative. Prove that for $C = O(k)$, with probability $9/10$, this scheme finds all truly positive patients and reports no false positives. Thus, we only require $O(k\log n)$ tests! 
	
%	\textbf{Hint:} If your proof would also work for $q = O(1)$ then it has a bug! See EC below.
	
	\item \textbf{Extra Credit -- optional.} (3pts) Show that no scheme can use $o(k\log(n/k))$ tests and succeed with probability $> 2/3$. So, for small $k$, the approach above is essentially optimal up to constant factors!
\end{enumerate}

%\subsection{Problem 4: Exploring Concentration Bounds, Part 1}
%\textbf{(8 pts)} Let $X$ be a random variable uniformly distributed in the interval $[0,1]$. Since we know $X$'s distribution exacty, we can easily check that $\Pr[X \geq 7/8] = 1/8$. But let's take a look at what various concentration inequalities would predict about this probability using less information about $X$. 
%\begin{enumerate}
%	\item Given an upper bound on $\Pr[X \geq 7/8]$ using Markov's inequality. 
%	\item Give an upper bound on $\Pr[X \geq 7/8]$ using Chebyshev's inequality.
%	\item Given an upper bound on $\Pr[X \geq 7/8]$ by applying Markov's inequality to the random variables $X^2$ (the ``raw'' second moment). 
%	Note that this is slightly different than using Chebyshev's inequality, which applies Markov to ``central'' second moment $(X - \E[X])^2$.
%	\item What happens for higher moments? Applying Markov's to $X^q$ for $q = 3,4, \ldots, 10$. Describe what you see in a sentence. What value of $q$ gives the tightest bound?
%		\item One take-away here is that, depending on the random variable being studied, it is not always optimal to use the variance as a deviation measure to show concentration. Markov's can be used with any monotonic function $g$, and as we see above, different choices might give better bounds. Exhibit a monotonic function $g$ so that applying Markov’s to $g(X)$ gives as tight an upper bound on $\Pr[X \geq 7/8]$ as you can. Maximum points if you can get $\Pr[X \geq 7/8] \leq 1/8$, which would be the best possible.
%\end{enumerate}
% f does not need to be strictly monotonic! Why? Sort of counterintuitive but got a question about this last year. 

% \subsection{Problem 4: Exploring Concentration Bounds. (8 pts)}
% Using combinatorics, write down and evaluate (e.g., run some code) an exact expression for the probability of flipping 60 or more heads when flipping a fair coin 100 times. Do the same for 600 or more heads in 1000 flips. Compare your results to the upper bound given by the Chernoff bound from Lecture 2:
% \begin{align*}
% 	\Pr[S > (1+\epsilon)\mu] \leq e^{-\frac{\epsilon^2\mu}{2+\epsilon}}.
% \end{align*}
% Also compare to this slightly different and less commonly used form of the bound from Wikipedia:
% \begin{align*}
% 	\Pr[S > (1+\epsilon)\mu] \leq \left(\frac{e^\epsilon}{(1+\epsilon)^{1+\epsilon}}\right)^{\mu}.
% \end{align*}
% Also compute what each Chernoff Bound would give if you wanted to bound the probability of flipping more than 70 or 700 heads. 
% Write a 1-2 sentence comment on how what your findings suggest about how ``tight'' the Chernoff bound is. 

\subsection{Problem 4: Try out mark-and-recapture! (12 pts)}
\begin{enumerate}
	\item (6pts) Prove the claim on Slide 31 of Lecture 1. Specifically, if we collection $O(\sqrt{n}/\epsilon)$ samples uniformly from a set of unknown size $n$, then we can return an estimation $\tilde{n}$ which, with probability $9/10$ satisfies:
	\begin{align*}
		(1-\epsilon)n \leq \tilde{n} \leq (1+\epsilon)n.
	\end{align*}
	
	\item (6pts) Wikipedia provides a way to access a random article by following the link \url{https://en.wikipedia.org/wiki/Special:Random}. Implement the mark-and-recapture algorithm from class and use this link to evaluate \href{https://en.wikipedia.org/wiki/English_Wikipedia}{the claim} that English Wikipedia has 6.7 million unique articles. In my experiments, downloading 5000 articles took 30 minutes, so scanning all possible articles to check the claim would take $\sim 25$ days (if you don't get blacklisted for scrapping first). How long did your code take to run?
	
	Include your code, your best estimte for the number of articles on Wikipedia, and how many samples you used to compute the estimate. 
	
	 I  used Python, but you can use any language. In Python, you can get a random url by running:
\begin{verbatim}
import requests
response = requests.get("https://en.wikipedia.org/wiki/Special:Random")
random_url = response.url	
\end{verbatim}

\item In part (2)My guess is that the underestimate is due to the fact that 
Wikipedia's random article generator does not return \emph{truly uniform} random articles. As discussed \href{https://en.wikipedia.org/wiki/Wikipedia:FAQ/Technical#random}{here}, Wikipedia assigns each article $i$ a random id $r_i$, which we can model as a random real number in $[0, 1]$. Then, to pick a random article, a random real number in $[0, 1]$ is sampled and article $i$ is return if that number lies in the range $[r_i, r_{i+1}].$ That is, the range $[0, 1]$ is partitioned into $n$ intervals and the probability of returning article $i$ is equal to the length of the $i^\text{th}$ interval, which we denote by $p_i$. Since these intervals themselves are random, the probability distribution won't be \emph{perfectly} uniform.
\begin{enumerate}
	\item Prove that, when the interval lengths are not perfectly uniform, the mark-and-recapture method should be expected to underestimate the number of articles. \textbf{Hint:} Prove that the expected number of collisions $D$ after $m$ samples equals $\binom{m}{2}\sum_{i=1}^n p_i^2$, and show that this expectation is minimized when $p_1 = \ldots = p_n = \frac{1}{n}$.  So, we get too many collisions for a non-uniform distribution, causing an underestimate of $n$.
	\item Show that, if Wikipedia uses the scheme above, we expect that the mark-and-recapture method will systematically underestimate the number of articles by almost exactly a factor of two. \textbf{Hint:} Try to prove a bound on $\E\left[\sum_{i=1}^n p_i^2\right]$. 
\end{enumerate}
\end{enumerate}


%\subsection{Problem 6: Randomized Routing}
%\textbf{(12 pts)}Randomization is often used in designing protocols for efficiently routing information over a network (e.g. the internet, or a network within a data center). Consider the following simplified model of a routing problem under communication constraints: we have $n$ nodes in a fully connected network (e.g. a data center with $n$ servers). Each node has $n$ packets that it would like to deliver to $n$ (not necessarily unique) destinations. Also, for simplicity, we assume that each node is the intended recipient of exactly $n$ packets. In each round, a node can send at most one packet along each connection in the network. Assume that the packet contains the id of its final intended recipient. 
%
%Naively, sending all packets could take $n$ rounds of communication, if e.g., some node $v$ wants to send all $n$ of its packets to some other node $u$.
%Design and analyze a randomized scheme that instead routes all packets to their intended recipient in $O(\log n)$ rounds with high probability, i.e., with probability at least $99/100$.
%\textbf{Hint:} Have each node initially send each of its $n$ packets to a random recipient.












\end{document}